% !TeX root = ../papers/nc-p3.tex
\section{長さ七以下の平面切断と二つの局所点の接平面}
\label{sec:finite-plane-sections-and-tangent-incidence}
\begin{theorem}\label{thm:finite-plane-section-seven}
四辺が有限な標準有限Frobenius窓の各座標において、点スキームの任意の平面切断は
長さ7以下である。長さ7なら、その平面と反対辺でない隣の三空間との積で、
一般の五関係が切る長さ10の完全交叉における残余は、平面へ共線な長さ3のスキームとして埋め込まれる。
\end{theorem}
\begin{proof}
隣接する辺を $\mathbb P^3\times\mathbb P^3$ の六双一次式とし、第一位置の平面 $H$ に制限する。
六式は独立なままである。依存すれば平面の式と第二位置の一次式との純関係が生じる。
基点は有限なので、一般の五式は $H\times\mathbb P^3$ に長さ $\binom52=10$ の完全交叉 $Y$ を与える。
$Y$ の $H$ への射影は閉埋込みである。全六式の評価が平面上の点で階数4なら五式でも少なくとも3、
階数3の有限個の点では一般の五式を選んでも階数3である。従って $Y$ はその退化スキーム上の核直線のグラフである。

残った一式の値 $f$ は $Y$ 上で零ではない。実際 $Y$ 上で消える双一次式は定義五式で尽きる。
Koszul分解を $(1,1)$ でひねると他の負の項は寄与しない（最終項の非零コホモロジーも次数5である）。
随伴により $\omega_Y=\mathcal O_Y(2,1)$ で、その完全跡は周囲から制限した $(2,1)$ 式を消す。
これはKoszul分解で単位写像に双対な最終接続写像がこの跡となるためである。
従って $f$ と $H$ の各一次式との積の跡は零。$Z=\operatorname{Spec}(\mathcal O_Y/\operatorname{ann} f)$ の
完全跡はこの $f$ 倍であり、平面の一次式は $Z$ に全射できない。平面射影が閉埋込みなので
$\operatorname{length}Z\ge3$ で、等号なら共線である。
元の切断長は $10-\operatorname{length} Z$ だから主張を得る。
\end{proof}

\begin{proposition}\label{prop:finite-two-planar-tangent-incidence}
反対階数8の有限窓で全中央支持4とする。二つの異なる平面局所点 $p,q$ について
$q_i\in T_{p_i}X_i$ は反対位置 $i+2$ での同じ包含と同値である。
一つの位置で $q_i\in T_{p_i}X_i$ と $p_i\in T_{q_i}X_i$ は両立しない。
また印 $p$ の両側のいずれかの分離三次式が $h s^2$ のとき、$X$ の台は一点ではあり得ない。
台が $p,q$ の二点で $q$ の局所長が8以上かつ $p$ の極大イデアルの二乗が非単項なら、四位置の全てで $q_i\in T_{p_i}X_i$ である。
\end{proposition}
\begin{proof}
偶側で $M_p=w(p_0,p_2)$ は階数1で、左右の因子は奇位置の接平面の法線である。
実際、印からの三座標に対して $M_p$ の左右核はその一次接方向を含み、各々二次元なので一致する。
反対切断の八空間の中で、$q_1$ による縮約が零の空間は次元4である。
それは $q$ の偶評価の共通四空間であり、$q_3$ での縮約が零である空間にも等しい。
従って $M_p$ の一方の法線が $q$ を消すことは他方でも同じである。添字を回せば偶位置側も従う。
両向きに包含すれば二点を結ぶ直線が両端で長さ2以上の切断を持つので、
系\ref{cor:finite-no-four-on-line}に反する。

二重因子の印付き切断を奇座標を変数として読む。定理\ref{thm:finite-split-separator-slices}により
$A=h vv^t+K$、$K$ は交代、$h(p)\ne0$ である。隣の偶位置の射影座標 $l$ に対して
$l^tA=0$ より $h(v^tl)^2=0$。$h$ が全ての局所因子で単元なら偶側の二次制限の単射性に反する。
従って一点支持は不可能で、二点支持なら $h(q)=0$。
$q$ の長さが8以上なら平面 $h=0$ にこの局所スキーム全体は入らない。
よって $v^tl$ は $q$ で非単元である。これは $p$ の偶接平面の式であり、反対位置でも同じである。
また分離式は $q$ の局所環で零、$h$ はそこで恒等零ではないので $s(q)=0$。
この $s$ の平面が奇変数での $p$ の接平面になることを、$p$ の局所環が非伸長型の場合に確かめる。
$E_D$ 表示の右下ブロックは $sD$、$a=h$ は単元、他の二座標を $b$ とすれば
局所式で $sD r'=a^{-1} b b^t r'$。$r'$ は二つの独立な一次方向を持つので
$s\mathfrak m\subset\mathfrak m^3$。もし $s$ の一次部が非零であれば
$\dim\mathfrak m^2/\mathfrak m^3\le1$ となる。従って非伸長型では $s\in\mathfrak m^2$ で所望の接平面である。
\end{proof}

\begin{corollary}\label{cor:two-support-quartic-discrepancy}
未認証の反対階数8の二点支持窓で、一つの印の階数一切断が長さ4とする。
核側が非対称のままなら、その外積の支持 $Z$ はその印だけに集中し、長さ3の曲線状環または
長さ4の局所 $(1,2,1)$ 環である。後者では $\Sigma$ も同じ一点に支持される。
どちらの場合もその印に重平面因子を持つ分離式が別の切断から得られる。
\end{corollary}
\begin{proof}
共線四点因子の排除と前節の結果から、残る支持環は上の二つまたは局所長 $2,1$ に分かれる長さ3である。
後者では単純因子での三次分離式は $H l_2^2$ である（$l_2$ は二重因子の方向）。
その局所点では $H$ が単元、$l_2$ は零であり、単純因子の方向 $l_1$ もその点で消える。
これは分離行列の階数1と外積支持の $Al=0$ による。従って $H,l_2,l_1$ は互いに比例せず、
二重因子の局所点の分離式 $H l_2 l_1$ は相異なる三平面の積になる。
定理\ref{thm:finite-split-separator-slices}からその印の切断は長さ2以下または正次元となり、既出の平面点判定で認証または排除できる。
最後に $Z$ が印を含まなければ、定理\ref{thm:finite-stretched-planar-marks}の
$\mathfrak m^2 f\subset(\eta)$ は印で $\mathfrak m^2 f=0$ を与える。
この包含自体の証明は伸長型を仮定しない。局所支持長は高々2なので、同定理の長さ下界との矛盾がそのまま成り立つ。
よって残る $Z$ は印そのものに支持され、分離式は $H s^2$ である。
\end{proof}
